\section*{第 \ref{ch:b3:asymmetric-synthesis} 章——不对称合成}

\begin{solution}{exo:b3:asymmetric-synthesis:1}
ee 为 80\,\%：$\mathrm{er} = 1.8/0.2 = 9$（90\,:\,10）。er 为 99\,:\,1 时：$\mathrm{ee} = 98/100 = 98$\,\%。ee 为 50\,\% 时：75\,\% 与 25\,\%。
\end{solution}

\begin{solution}{exo:b3:asymmetric-synthesis:2}
$(1840 - 60)/(1840 + 60) = 0.937$：ee 为 93.7\,\%，er 为 $1840/60 = 31$。
\end{solution}

\begin{solution}{exo:b3:asymmetric-synthesis:3}
优先级 O $>$ 乙基 $>$ H。以 O 朝上、乙基在左下、H 在右下画出时，
前面为 \textit{Si} 面，后面为 \textit{Re} 面。从 Re 面加入甲基，
产物中优先级为 OH $>$ 乙基 $>$ 甲基 $>$ H，得到（\textit{R}）-2-丁醇。
\end{solution}

\begin{solution}{exo:b3:asymmetric-synthesis:4}
乙醇：对映异位（分子的\omterm{def:b3:point-groups:mirror}{镜面}可使其互换；
在非手性溶剂中化学位移相同。2-丁醇：非对映异位（与手性碳原子相邻；
化学位移不同）。
\end{solution}

\begin{solution}{exo:b3:asymmetric-synthesis:5}
er 为 199：$RT\ln 199 = \qty{13.1}{kJ/mol}$（\qty{298}{K}），\qty{8.6}{kJ/mol}（\qty{195}{K}）。\qty{25}{\celsius} 下 80\,\% ee 对应 er 为 9，
$\Delta\Delta G^\ddagger = RT\ln 9 = \qty{5.45}{kJ/mol}$；在 \qty{-78}{\celsius} 下，$\mathrm{er} = \eu^{5450/(8.314 \times 195.15)} = 28.7$，ee 为 93\,\%。
\end{solution}

\begin{solution}{exo:b3:asymmetric-synthesis:6}
\omterm{def:b3:asymmetric-synthesis:ee}{光学纯度}为 $12.0/15.0 = 80$\,\%：90\,\%（\textit{S}）与 10\,\%（\textit{R}）。旋光度较小、存在杂质、
浓度与溶剂效应，以及旋光度与组成之间的非线性关系，
都使其不可靠，色谱法可直接分离并计量
对映异构体。
\end{solution}

\begin{solution}{exo:b3:asymmetric-synthesis:7}
$s = \ln(0.40 \times 0.10)/\ln(0.40 \times 1.90) = -3.22/-0.274 = 12$.
\end{solution}

\begin{solution}{exo:b3:asymmetric-synthesis:8}
使用 L-($+$)-酒石酸二乙酯时为（2\textit{S},3\textit{S}）-2,3-环氧己-1-醇；
使用 D-($-$)-酒石酸二乙酯时，得其对映异构体（2\textit{R},3\textit{R}）。
\end{solution}

\begin{solution}{exo:b3:asymmetric-synthesis:9}
L 为苯基，M 为甲基，S 为 H。对于（\textit{S}）醛，反应构象中苯基
垂直于羰基平面，H 位于亲核试剂进攻路径旁，
亲核试剂从苯基的反面进攻。由此推出构型，
主要非对映异构体为（2\textit{S},3\textit{S}）-3-苯基-2-丁醇。
\end{solution}

\begin{solution}{exo:b3:asymmetric-synthesis:10}
反应开始时两种对映异构体的量相等，因此产物之比为
$k_{\mathrm{fast}} : k_{\mathrm{slow}} = s$：$\mathrm{ee}_p = (s - 1)/(s + 1)$，$s = 20$ 时为 0.905。
随着反应进行，快速对映异构体被消耗，产物 ee 下降，
而底物 ee 趋向 1：底物总可通过继续反应得到提纯，产物则
不能。
\end{solution}

\begin{solution}{exo:b3:asymmetric-synthesis:11}
快速外消旋化时：产率可达 100\,\%，$\mathrm{er} = 99$，ee 为 98\,\%。无外消旋化时，
转化率为 50\,\% 时，产率至多为 50\,\%，产物 ee 为 93\,\%；
底物的 ee 同样为 93\,\%。
\end{solution}

\begin{solution}{exo:b3:asymmetric-synthesis:12}
两种配合物互变速率快于反应速率时，产物之比为
$k_B[\mathrm B]/(k_A[\mathrm A]) = 600/10 = 60$：次要配合物 B 生成 98\,\% 的产物（ee 为 97\,\%）。
选择性由两个过渡态的能量差决定，而非
由中间体的\omterm{def:b3:partition-functions:boltzmann}{布居数}决定。
\end{solution}

\begin{solution}{pb:b3:asymmetric-synthesis:1}
\textbf{1.} 带有氮的烯烃碳连有三个不同的取代基，
并具有两个不同的面；向其上加氢后，
氨基酸的 $\alpha$ 碳即成为一个手性碳原子。

\textbf{2.} 外消旋混合物：两个面被进攻的速率相同。

\textbf{3.} 其羰基氧与铑及 C=C 键同时配位，底物以
螯合物形式结合，几何构型固定，配体可对其加以区分。

\textbf{4.} 它在金属周围构造一个手性口袋（$C_2$ 对称）：底物
从一面或另一面结合的两种方式互为非对映异构体，具有
不同的能量与反应活性。

\textbf{5.} 不对：次要配合物与氢的反应要快得多（柯廷–哈米特）。

\textbf{6.} 金属位于已配位烯烃的一面，
氢原子从金属转移而来。

\textbf{7.} $\mathrm{er} = 1.95/0.05 = 39$.

\textbf{8.} $\Delta\Delta G^\ddagger = RT\ln 39 = 8.314 \times 298.15 \times 3.664 = \qty{9.1}{kJ/mol}$.

\textbf{9.} $\mathrm{er} = \eu^{9082/(8.314 \times 195.15)} = 270$: ee 99.3\,\%.

\textbf{10.} $\mathrm{er} = \eu^{11082/(8.314 \times 298.15)} = 87$: ee 97.7\,\%.

\textbf{11.} 比典型氢键更弱：底物与配体之间空间接触的细微差异，
或一个弱相互作用，即可决定结果，
这也是此类催化剂难以从第一性原理出发进行设计的原因。

\textbf{12.} 反应速率下降，氢气的溶解情况改变，
催化剂在低温下的表现也不同，冷却大型反应器还需耗费能量。

\textbf{13.} （\textit{S}）对映异构体 97.5 g，（\textit{R}）对映异构体 2.5 g。

\textbf{14.} $97.5 - 5.0 = \qty{92.5}{g}$ 晶体，实际为对映纯的。

\textbf{15.} $92.5/100 = 92.5$\,\%（以粗产物计，其中（\textit{S}）对映异构体占 95\,\%）。

\textbf{16.} 晶体吸收过量的对映异构体，
外消旋部分与少量次要对映异构体留在溶液中；
从外消旋混合物中无法收集到过量部分，结晶得到的是外消旋晶体。

\textbf{17.} 用手性 HPLC，以外消旋标样作参照。

\textbf{18.} 50\,\%：只有慢速对映异构体保留下来。

\textbf{19.} $s = \ln(0.45 \times 0.05)/\ln(0.45 \times 1.95) = -3.79/-0.131 = 29$.

\textbf{20.} 外消旋混合物的 45\,\%（起始时所需对映异构体的 88\,\%，
ee 为 95\,\%）。

\textbf{21.} 若底物快速外消旋化，产率可接近 100\,\%。

\textbf{22.} 它将全部廉价的非手性前体转化为正确的对映异构体，
所用催化剂量极少，且无需丢弃任何对映异构体。

\textbf{23.} 铑有毒，药物中其残留量限于百万分之几；
其含量采用电感耦合等离子体质谱法测定。

\textbf{24.} \qty{25}{\celsius} 下 95\,\% ee 对应 $\Delta\Delta G^\ddagger = RT\ln 39 \approx \qty{9.1}{kJ/mol}$。
\end{solution}
